%==============================================================================
%
%  APPENDIX V: THE 8Ψ-DIMENSIONS FRAMEWORK
%  Context Assessment for Behavioral Economic Analysis
%
%  EBF CORE Appendix - Answers: WHEN does context matter?
%
%==============================================================================

%------------------------------------------------------------------------------
% 1.1 HEADER BLOCK (Metadata)
%------------------------------------------------------------------------------

% Category: LIT
% Language: English
\begin{tcolorbox}[colback=green!5!white, colframe=green!75!black, 
    title=\textbf{CORE Appendix: Context Assessment}]

\begin{tabular}{@{}ll@{}}
\textbf{Appendix:} & V (8$\Psi$-Dimensions Framework) \\
\textbf{Category:} & CORE \\
\textbf{CORE Question:} & \textbf{WHEN} --- When does context matter? \\
\textbf{Scope:} & A3 (Methodological Framework) \\
\textbf{Version:} & 2.0 \\
\textbf{Last Updated:} & January 2026 \\
\textbf{Dependencies:} & Appendix HOW (Complementarity Theory) \\
 & Appendix WAT (FEPSDE Utility Structure) \\
 & Appendix EST (Parameter Calibration) \\
\textbf{Outputs:} & Context vector $\boldsymbol{\Psi} \in [0,1]^8$ \\
 & $\delta$-modulation coefficients \\
 & Context-adjusted complementarities $\gamma(\Psi)$ \\
\end{tabular}

\end{tcolorbox}

%------------------------------------------------------------------------------
% 1.2 CHAPTER LINKAGE BOX
%------------------------------------------------------------------------------

\begin{tcolorbox}[colback=gray!10!white, colframe=gray!75!black, 
    title=\textbf{Chapter Linkages}]

This appendix provides methodological foundations for:

\begin{itemize}[nosep]
    \item \textbf{Chapter 3} (Model Specification): Context vector $\boldsymbol{\Psi}$ enters utility functions
    \item \textbf{Chapter 5} (Interventions): Context determines intervention effectiveness
    \item \textbf{Chapter 7} (Predictions): Treatment effect heterogeneity across contexts
    \item \textbf{Chapter 8} (Applications): Country/region-specific calibration
\end{itemize}

\textbf{Key insight:} The same intervention produces different effects in different 
contexts. The 8$\Psi$ framework explains \textit{why} and predicts \textit{how much}.

\end{tcolorbox}

%------------------------------------------------------------------------------
% 1.3 CROSS-REFERENCE MAP
%------------------------------------------------------------------------------

\begin{tcolorbox}[colback=white, colframe=black, 
    title=\textbf{Cross-Reference Map}]

\begin{center}
\small
\begin{tabular}{@{}c|ccccc@{}}
\toprule
& \textbf{B} & \textbf{C} & \textbf{AAA} & \textbf{BBB} & \textbf{AN} \\
\midrule
\textbf{V $\leftarrow$} & $\gamma$ theory & FEPSDE & Levels & Calibration & LLM-MC \\
\textbf{V $\rightarrow$} & $\gamma(\Psi)$ & $w_d(\Psi)$ & $\Psi$ by level & $\delta$-coefficients & Validation \\
\midrule
\textbf{Link} & \cellcolor{red!20}Strong & \cellcolor{red!20}Strong & \cellcolor{yellow!20}Medium & \cellcolor{red!20}Strong & \cellcolor{yellow!20}Medium \\
\bottomrule
\end{tabular}
\end{center}

\textbf{Legend:}
\colorbox{red!20}{Strong} = Bidirectional dependency \quad
\colorbox{yellow!20}{Medium} = One-way dependency

\textbf{Additional Reference:}
\begin{itemize}[nosep]
    \item \textbf{XVIII (LIT-FEHR-METHOD):} Exclusion Principle --- $\gamma = 0$ as null hypothesis
\end{itemize}

\end{tcolorbox}

%==============================================================================
% FORMAL COMPLIANCE BLOCK
%==============================================================================

% --- Terminology Reference ---
\begin{tcolorbox}[colback=gray!5!white, colframe=gray!50!black]
\textbf{Terminology:} See Appendix GLS (Glossary) for standard definitions.
Key terms in this appendix: Context vector ($\boldsymbol{\Psi} \in [0,1]^8$),
level salience $\alpha^L(\Psi)$, context modulation $\delta(\Psi)$,
8$\Psi$ dimensions ($\Psi_I, \Psi_S, \Psi_C, \Psi_K, \Psi_E, \Psi_T, \Psi_M, \Psi_F$).
\end{tcolorbox}

% --- Epistemic Status Tags ---
\begin{tcolorbox}[colback=blue!5!white, colframe=blue!75!black,
    title=\textbf{Epistemic Status: Key Parameters}]

\renewcommand{\arraystretch}{1.3}
\begin{tabular}{@{}llcc@{}}
\toprule
\textbf{Parameter} & \textbf{Description} & \textbf{Proxy/Range} & \textbf{Tag} \\
\midrule
$\Psi_I$ & Institutional quality & Rule of Law Index & [EMP] \\
$\Psi_S$ & Social capital & World Values Survey & [EMP] \\
$\Psi_C$ & Cognitive capacity & PISA scores & [EMP] \\
$\Psi_K$ & Information quality & Press Freedom Index & [EMP] \\
$\Psi_E$ & Economic development & GDP per capita (log) & [EMP] \\
$\Psi_T$ & Temporal stability & Policy volatility & [EMP] \\
$\Psi_M$ & Market scope & Trade openness & [EMP] \\
$\Psi_F$ & Factor flexibility & Labor market rigidity & [EMP] \\
\addlinespace
\multicolumn{4}{@{}l}{\textit{Modulation effects:}} \\
$\delta_I$ & Institutional modulation & 0.1--0.4 & [LLM] \\
$\delta_S$ & Social modulation & 0.1--0.3 & [LLM] \\
$R^2$ & Heterogeneity explained & $\approx 55\%$ & [EMP] \\
\bottomrule
\end{tabular}

\vspace{0.5em}
\textbf{Tag Legend:} [EMP] = Empirically validated (cross-country data);
[LLM] = LLM Monte Carlo estimated

\end{tcolorbox}

% --- Bibliography Reference ---
\noindent\textbf{Full citations:} See \texttt{bibliography/bcm\_master.bib} \\
\textbf{Key sources:} North (1990) institutions; Putnam (1993) social capital; World Bank Governance Indicators

\vspace{1em}

%------------------------------------------------------------------------------
% 1.4 ABSTRACT
%------------------------------------------------------------------------------

\begin{tcolorbox}[colback=white, colframe=black,
    title=\textbf{Abstract}]

Economic effects are context-dependent, yet ``context'' often serves as an 
unexplained residual. This appendix operationalizes context through eight 
primitive dimensions ($\Psi_I$ Institutional, $\Psi_S$ Social, $\Psi_C$ Cognitive, 
$\Psi_K$ Informational, $\Psi_E$ Economic, $\Psi_T$ Temporal, $\Psi_M$ Market Scope, 
$\Psi_F$ Factor Flexibility) derived from decision-theoretic first principles. 
Empirical validation across six canonical economic relationships shows the 
8$\Psi$ model explains 55\% of treatment effect heterogeneity on average. 
The framework enables systematic context assessment for intervention design 
and cross-context prediction.

\textit{(98 words)}

\end{tcolorbox}

%------------------------------------------------------------------------------
% 1.5 QUICK REFERENCE BOX
%------------------------------------------------------------------------------

\begin{tcolorbox}[colback=yellow!10!white, colframe=yellow!75!black, 
    title=\textbf{Quick Reference}]

\textbf{The Context Vector:}
\[
\boldsymbol{\Psi} = (\Psi_I, \Psi_S, \Psi_C, \Psi_K, \Psi_E, \Psi_T, \Psi_M, \Psi_F) \in [0,1]^8
\]

\textbf{The Eight Dimensions:}
\begin{center}
\small
\begin{tabular}{@{}clll@{}}
\toprule
\textbf{Symbol} & \textbf{Dimension} & \textbf{Measures} & \textbf{High $\Psi$ means...} \\
\midrule
$\Psi_I$ & Institutional & Rule of Law & Strong enforcement \\
$\Psi_S$ & Social & Generalized Trust & High cooperation \\
$\Psi_C$ & Cognitive & PISA scores & High processing capacity \\
$\Psi_K$ & Informational & Press Freedom & Transparent information \\
$\Psi_E$ & Economic & GDP/capita & High development \\
$\Psi_T$ & Temporal & Policy Stability & Long time horizons \\
$\Psi_M$ & Market Scope & Trade openness & Global integration \\
$\Psi_F$ & Factor Flexibility & Labor mobility & Easy adjustment \\
\bottomrule
\end{tabular}
\end{center}

\textbf{Context Modulation:}
\[
\gamma_{dd'}(\boldsymbol{\Psi}) = \bar{\gamma}_{dd'} + \sum_{k=1}^{8} \delta_{dd',k} \cdot (\Psi_k - 0.5)
\]

\textbf{Jump to:}
\begin{itemize}[nosep]
    \item Theoretical derivation $\rightarrow$ \S\ref{sec:V-theory}
    \item Empirical validation $\rightarrow$ \S\ref{sec:V-validation}
    \item Measurement protocol $\rightarrow$ \S\ref{sec:V-measurement}
    \item Worked example (Switzerland) $\rightarrow$ \S\ref{sec:V-example}
    \item Axiom system $\rightarrow$ \S\ref{sec:V-axioms}
\end{itemize}

\end{tcolorbox}



% -----------------------------------------------------------------------------
% APPENDIX SCOPE BOX
% -----------------------------------------------------------------------------

\begin{tcolorbox}[
    colback=orange!5!white,
    colframe=orange!75!black,
    title={\textbf{Appendix Scope: Ziel / In-Scope / Out-of-Scope / Constraints}},
    fonttitle=\bfseries
]

\textbf{Ziel (Objective):}
\begin{itemize}[nosep]
    \item [TODO: Define objective for Unknown Title]
\end{itemize}

\textbf{In-Scope:}
\begin{itemize}[nosep]
    \item The Eight $\Psi$ Dimensions: Operationalization and Empirical Validation
    \item Axiom System for Context Dimensions
    \item Measurement Protocol
    \item Worked Example: Context Assessment for Switzerland
    \item Worked Example 2: Organizational Context Design
\end{itemize}

\textbf{Out-of-Scope (delegiert an andere Appendices):}
\begin{itemize}[nosep]
    \item Complementarity Theory $\rightarrow$ Appendix HOW
    \item FEPSDE Utility Structure $\rightarrow$ Appendix WAT
    \item Parameter Calibration $\rightarrow$ Appendix EST
    \item Glossary $\rightarrow$ Appendix GLS
    \item \S\ref{sec:bbb-delta} $\rightarrow$ Appendix EST
\end{itemize}

\textbf{Constraints (Anwendungsgrenzen):}
\begin{itemize}[nosep]
    \item [TODO: Define when this appendix does NOT apply]
    \item [TODO: Define required prerequisites]
    \item [TODO: Define known limitations]
\end{itemize}

\textbf{Lieferobjekte (Deliverables):}
\begin{itemize}[nosep]
    \item 11 Table(s)
    \item 14 Equation(s)
    \item Worked Example(s)
\end{itemize}

\end{tcolorbox}


\section{The Eight $\Psi$ Dimensions: Operationalization and Empirical Validation}
\label{app:psi-dimensions}

This appendix addresses a fundamental critique: that the context parameter $\Psi$ is merely a residual category---a placeholder for ``everything that matters'' without theoretical content. We demonstrate that this critique is unfounded. The context space $\Psi$ can be decomposed into exactly eight primitive dimensions, each independently measurable and collectively sufficient to explain observed heterogeneity in economic effects across multiple domains.

\subsection{Theoretical Derivation}
\label{sec:V-theory}

We begin with a fundamental question: What must any economic agent know, or be situated within, to make a decision? Any decision requires:
\begin{itemize}
    \item \textbf{Who decides?} (Agent identity and characteristics)
    \item \textbf{What are the options?} (Choice set)
    \item \textbf{What are the consequences?} (Payoff structure)
    \item \textbf{What rules apply?} (Constraints and enforcement)
    \item \textbf{What does the agent know?} (Information)
    \item \textbf{When must the decision be made?} (Temporal structure)
\end{itemize}

These requirements are logical necessities. From them, we derive eight context dimensions:
\begin{equation}
\Psi = (\Psi_I, \Psi_S, \Psi_C, \Psi_K, \Psi_E, \Psi_T, \Psi_M, \Psi_F) \in \mathbb{R}^8
\label{eq:psi-vector-appv}
\end{equation}

\begin{table}[htbp]
\centering
\caption{The Eight $\Psi$ Dimensions: Definition and Operationalization}
\label{tab:psi-dimensions-appv}
\small
\begin{tabular}{@{}cllll@{}}
\toprule
\textbf{Symbol} & \textbf{Dimension} & \textbf{Definition} & \textbf{Measures} & \textbf{Sources} \\
\midrule
$\Psi_I$ & Institutional & Rules, enforcement (de jure) & Rule of Law, Regulatory Quality & WGI, Polity IV \\
$\Psi_S$ & Social & Trust, norms, networks & Generalized Trust & WVS, ESS \\
$\Psi_C$ & Cognitive & Processing capacity, biases & Cognitive Skills & PISA, PIAAC \\
$\Psi_K$ & Informational & Transparency, signal quality & Information Access & TI, Press Freedom \\
$\Psi_E$ & Economic & Market structure, resources & GDP/capita, Development & World Bank \\
$\Psi_T$ & Temporal & Stability, time horizons & Policy Uncertainty & EPU Index \\
$\Psi_M$ & Market Scope & Geographic reach, openness & Trade/GDP, FDI & KOF Index \\
$\Psi_F$ & Factor Flexibility & Adjustment capacity & Labor Market Flexibility & OECD EPL \\
\bottomrule
\end{tabular}
\end{table}

\subsection{Empirical Validation: Six Tests}
\label{sec:V-validation}

We test whether these eight dimensions explain treatment effect heterogeneity across six canonical economic relationships.

\subsubsection{Test 1: Education $\rightarrow$ Earnings}

Returns to education vary from 5.5\% in Scandinavia to over 14\% in parts of Africa. 

\textbf{Result:} Adjusted $R^2 = 0.768$. Dominant dimensions: $\Psi_K$ (information) and $\Psi_M$ (market scope). Education returns depend on signaling value and access to skill-rewarding labor markets.

\subsubsection{Test 2: Management Practices $\rightarrow$ Productivity}

Following \citet{bloom2012management}, management practices affect productivity differently across countries.

\textbf{Result:} Adjusted $R^2 = 0.683$. Dominant dimension: $\Psi_F$ (factor flexibility). Management practices only improve productivity when firms can implement changes---hiring, firing, reorganizing. In rigid labor markets, even excellent management knowledge cannot translate into productivity gains.

\subsubsection{Test 3: Democracy $\rightarrow$ Growth}

The effect of democracy on growth varies widely.

\textbf{Result:} Adjusted $R^2 = 0.097$. Dominant dimensions: $\Psi_I$, $\Psi_E$. Macro-level political effects are inherently noisier than micro-level decisions.

\subsubsection{Test 4: Patience $\rightarrow$ Growth}

Using \citet{PAP-enke2019globalpref} Global Preference Survey data, we test whether patience produces different growth effects depending on context.

\textbf{Result:} Adjusted $R^2 = 0.510$. Dominant dimensions: $\Psi_T$ (temporal stability), $\Psi_K$ (information). Patience provides greater advantage in uncertain, opaque environments.

\subsubsection{Test 5: Information $\rightarrow$ Behavior Change}

Information interventions show highly heterogeneous effects \citep{jensen2007digital, dupas2011teenagers}.

\textbf{Result:} Adjusted $R^2 = 0.737$. Dominant dimensions: $\Psi_K$ (\textit{negatively}), $\Psi_F$. 

\textbf{Critical finding:} Information interventions work when (1) $\Psi_K$ is \textit{low}---the information is genuinely new, and (2) $\Psi_F$ is \textit{high}---people can act on it. This explains why \citet{jensen2007digital} found large effects (fishermen could immediately change selling location) while \citet{dupas2011teenagers} found small effects (HIV information was already known; social constraints prevented behavior change).

\subsubsection{Test 6: Income $\rightarrow$ Life Expectancy}

The Preston Curve shows richer countries live longer, but with substantial residual variance.

\textbf{Result:} Adjusted $R^2 = 0.516$. Dominant dimensions: $\Psi_T$ (stability), $\Psi_I$ (institutions). Countries with peace and functioning institutions live longer than income alone predicts.

\subsection{Summary of Results}

\begin{table}[htbp]
\centering
\caption{Empirical Validation: 8$\Psi$ Model Across Six Outcomes}
\label{tab:psi-validation-appv}
\small
\begin{tabular}{@{}lcccc@{}}
\toprule
\textbf{Economic Relationship} & \textbf{Type} & \textbf{$R^2$} & \textbf{Adj. $R^2$} & \textbf{Dominant $\Psi$} \\
\midrule
Education $\rightarrow$ Earnings & Micro & 84.5\% & 76.8\% & $\Psi_K$, $\Psi_M$ \\
Management $\rightarrow$ Productivity & Micro & 78.9\% & 68.3\% & $\Psi_F$ \\
Democracy $\rightarrow$ Growth & Macro & 39.8\% & 9.7\% & $\Psi_I$, $\Psi_E$ \\
Patience $\rightarrow$ Growth & Macro & 67.3\% & 51.0\% & $\Psi_T$, $\Psi_K$ \\
Information $\rightarrow$ Behavior & Behavioral & 82.4\% & 73.7\% & $\Psi_K$ (--), $\Psi_F$ \\
Income $\rightarrow$ Life Expectancy & Health & 67.7\% & 51.6\% & $\Psi_T$, $\Psi_I$ \\
\midrule
\textbf{Average} & --- & \textbf{70.1\%} & \textbf{55.2\%} & --- \\
\bottomrule
\end{tabular}
\end{table}

\subsection{Key Findings}

Three findings emerge from the validation:

\begin{enumerate}
    \item \textbf{Different dimensions dominate for different effects.} $\Psi_F$ matters most for management; $\Psi_K$ (negatively) for information interventions; $\Psi_T$ for patience effects. Same context space, different weights.
    
    \item \textbf{Micro effects are better explained than macro effects.} Individual and firm-level decisions ($R^2 \approx 70$--$85\%$) show more systematic context-dependence than aggregate outcomes ($R^2 \approx 40$--$70\%$). This is expected: aggregation introduces noise.
    
    \item \textbf{No ninth dimension emerged.} We tested candidates including implementation efficiency ($\Psi_X$), ethnic fractionalization, and natural resources. None significantly improved the model across all outcomes.
\end{enumerate}



%==============================================================================
% AXIOM SYSTEM FOR THE 8Ψ-DIMENSIONS
%==============================================================================

\section{Axiom System for Context Dimensions}
\label{sec:V-axioms}

The 8$\Psi$-Dimensions are not arbitrary---they are derived from axiomatic 
foundations. This section formalizes the properties that any valid context 
framework must satisfy.

\subsection{Foundational Axioms}

\begin{axiom}[Ψ-EXIST: Existence of Context]
\label{ax:psi-exist}
For any economic decision $a$ by agent $i$, there exists a context vector 
$\boldsymbol{\Psi} \in \mathcal{C}$ such that the utility outcome depends on both:
\[
U_i(a) = U_i(a; \boldsymbol{\Psi})
\]
Context is not eliminable from utility analysis.
\end{axiom}

\begin{axiom}[Ψ-DIM: Finite Dimensionality]
\label{ax:psi-dim}
The context space $\mathcal{C}$ can be represented by a finite-dimensional 
vector space:
\[
\mathcal{C} \cong [0,1]^K \quad \text{for some } K < \infty
\]
We claim $K = 8$ is necessary and sufficient (see Axiom \ref{ax:psi-complete}).
\end{axiom}

\begin{axiom}[Ψ-BOUND: Bounded Context]
\label{ax:psi-bound}
Each context dimension is bounded and normalized:
\[
\Psi_k \in [0,1] \quad \forall k \in \{I, S, C, K, E, T, M, F\}
\]
where $\Psi_k = 0$ represents the theoretical minimum and $\Psi_k = 1$ 
the theoretical maximum of dimension $k$.
\end{axiom}

\subsection{Structural Axioms}

\begin{axiom}[Ψ-COMPLETE: Completeness]
\label{ax:psi-complete}
The eight dimensions are \textbf{jointly sufficient} to characterize any 
economically relevant context:
\[
\forall \text{ treatment effect } \tau: \quad 
\text{Var}(\tau | \boldsymbol{\Psi}) \ll \text{Var}(\tau)
\]
Adding a ninth dimension does not systematically improve explanatory power 
across all outcome domains.

\textbf{Empirical support:} We tested candidates including ethnic fractionalization, 
natural resources, and implementation capacity. None significantly improved 
the model across all six validation tests (\S\ref{sec:V-validation}).
\end{axiom}

\begin{axiom}[Ψ-NECESSARY: Necessity]
\label{ax:psi-necessary}
Each of the eight dimensions is \textbf{necessary}---removing any dimension 
significantly reduces explanatory power for at least one class of economic effects:
\[
\forall k: \exists \text{ outcome } Y_k \text{ such that } 
R^2(Y_k | \boldsymbol{\Psi}) \gg R^2(Y_k | \boldsymbol{\Psi}_{-k})
\]
where $\boldsymbol{\Psi}_{-k}$ denotes the context vector with dimension $k$ removed.
\end{axiom}

\begin{axiom}[Ψ-ORTH: Approximate Orthogonality]
\label{ax:psi-orth}
The dimensions are \textbf{conceptually distinct} and \textbf{empirically 
separable}, though not perfectly orthogonal:
\[
|\text{Corr}(\Psi_k, \Psi_l)| < 0.7 \quad \forall k \neq l
\]
Moderate correlations are expected (e.g., institutions and development), 
but each dimension captures unique variance.
\end{axiom}

\begin{axiom}[Ψ-MEAS: Measurability]
\label{ax:psi-meas}
Each dimension $\Psi_k$ is \textbf{operationalizable} through observable 
indicators:
\[
\Psi_k = f_k(\text{observable data})
\]
with measurement error $\epsilon_k$ satisfying $\mathbb{E}[\epsilon_k] = 0$ 
and $\text{Var}(\epsilon_k) < \sigma^2_{\max}$.

Multiple proxy measures exist for each dimension (see \S\ref{sec:V-measurement}).
\end{axiom}

\subsection{Dynamic Axioms}

\begin{axiom}[Ψ-STABLE: Conditional Stability]
\label{ax:psi-stable}
Context dimensions change slowly relative to individual decisions:
\[
\frac{\partial \Psi_k}{\partial t} \ll \frac{\partial a_i}{\partial t}
\]
For most analyses, $\boldsymbol{\Psi}$ can be treated as fixed within a 
decision period. Exceptions: crises, regime changes.
\end{axiom}

\begin{axiom}[Ψ-ENDO: Potential Endogeneity]
\label{ax:psi-endo}
Context is \textbf{potentially endogenous} to aggregate behavior:
\[
\boldsymbol{\Psi}_{t+1} = g(\boldsymbol{\Psi}_t, \{a_i\}_{i \in \mathcal{N}})
\]
Institutions, trust, and norms evolve based on collective action. 
This creates feedback loops requiring careful identification.

\textbf{Implication:} Cross-sectional analysis treats $\boldsymbol{\Psi}$ 
as exogenous; dynamic analysis must model $\boldsymbol{\Psi}$ evolution.
\end{axiom}

\subsection{Modulation Axioms}

\begin{axiom}[Ψ-MOD: Linear Modulation]
\label{ax:psi-mod}
Context modulates complementarity through an approximately linear relationship:
\[
\gamma_{dd'}(\boldsymbol{\Psi}) = \bar{\gamma}_{dd'} + \sum_{k=1}^{8} 
\delta_{dd',k} \cdot (\Psi_k - 0.5) + O(\Psi^2)
\]
The $\delta$-coefficients capture the marginal effect of context on 
complementarity. Second-order terms are negligible for most applications.
\end{axiom}

\begin{axiom}[Ψ-HETERO: Heterogeneous Effects]
\label{ax:psi-hetero}
Different utility pairs $(d, d')$ are modulated by different context dimensions:
\[
\delta_{dd',k} \neq \delta_{d'd'',k} \quad \text{in general}
\]
There is no universal ``most important'' context dimension---relevance 
depends on the specific complementarity being analyzed.

\textbf{Example:} $\Psi_F$ (flexibility) dominates for management-productivity 
effects; $\Psi_K$ (information) dominates for education-earnings effects.
\end{axiom}

\subsection{Axiom Summary}

\begin{table}[htbp]
\centering
\caption{Complete Axiom System for 8$\Psi$-Dimensions}
\label{tab:V-axiom-summary}
\small
\renewcommand{\arraystretch}{1.2}
\begin{tabular}{@{}llp{7cm}@{}}
\toprule
\textbf{Axiom} & \textbf{Type} & \textbf{Statement} \\
\midrule
Ψ-EXIST & Foundational & Context affects utility outcomes \\
Ψ-DIM & Foundational & Context is finite-dimensional \\
Ψ-BOUND & Foundational & Dimensions are bounded on $[0,1]$ \\
\midrule
Ψ-COMPLETE & Structural & 8 dimensions are jointly sufficient \\
Ψ-NECESSARY & Structural & Each dimension is individually necessary \\
Ψ-ORTH & Structural & Dimensions are approximately orthogonal \\
Ψ-MEAS & Structural & Each dimension is measurable \\
\midrule
Ψ-STABLE & Dynamic & Context changes slowly \\
Ψ-ENDO & Dynamic & Context is potentially endogenous \\
\midrule
Ψ-MOD & Modulation & Linear modulation of complementarity \\
Ψ-HETERO & Modulation & Heterogeneous effects across pairs \\
\bottomrule
\end{tabular}
\end{table}



\subsection{Detailed Data Sources and Methodology}
\label{sec:psi-data-details}

This section provides complete documentation for replicating the empirical validation. For each test, we specify: (1) the source of treatment effect estimates, (2) the sample of countries, (3) the exact $\Psi$ measures used, and (4) the regression specification.

\subsubsection{Test 1: Education $\rightarrow$ Earnings}

\paragraph{Treatment Effect Source.}
Returns to education (Mincerian coefficients) from \citet{psacharopoulos2018returns}: ``Returns to Investment in Education: A Decennial Review of the Global Literature,'' \textit{Education Economics}, 26(5), 445--458. This meta-analysis provides comparable estimates across countries using consistent methodology.

\paragraph{Sample.}
$N = 28$ countries with complete data: Argentina, Australia, Austria, Belgium, Brazil, Canada, Chile, China, Colombia, Denmark, Finland, France, Germany, Greece, India, Ireland, Italy, Japan, Mexico, Netherlands, Norway, Poland, South Korea, Spain, Sweden, Switzerland, United Kingdom, United States.

\paragraph{Outcome Variable.}
$\tau_i$ = Private returns to one additional year of schooling (\%), averaged 2010--2017.

\paragraph{$\Psi$ Measures (all averaged 2010--2015).}
\begin{center}
\small
\begin{tabular}{@{}lll@{}}
\toprule
\textbf{Dimension} & \textbf{Measure} & \textbf{Source} \\
\midrule
$\Psi_I$ & Rule of Law Index (--2.5 to +2.5) & World Bank WGI \\
$\Psi_S$ & ``Most people can be trusted'' (\%) & World Values Survey Wave 6 \\
$\Psi_C$ & PISA Mathematics Score (mean) & OECD PISA 2012/2015 \\
$\Psi_K$ & Press Freedom Index (inverted, 0--100) & Reporters Without Borders \\
$\Psi_E$ & ln(GDP per capita, PPP, constant 2017\$) & World Bank WDI \\
$\Psi_T$ & Political Stability Index (--2.5 to +2.5) & World Bank WGI \\
$\Psi_M$ & Trade Openness (Exports+Imports)/GDP & World Bank WDI \\
$\Psi_F$ & Employment Protection Legislation (inverted) & OECD EPL Index \\
\bottomrule
\end{tabular}
\end{center}

\paragraph{Specification.}
\begin{equation}
\tau_i^{educ} = \beta_0 + \beta_I \Psi_{I,i} + \beta_S \Psi_{S,i} + \beta_C \Psi_{C,i} + \beta_K \Psi_{K,i} + \beta_E \Psi_{E,i} + \beta_T \Psi_{T,i} + \beta_M \Psi_{M,i} + \beta_F \Psi_{F,i} + \varepsilon_i
\end{equation}

\paragraph{Results.}
$R^2 = 0.845$, Adjusted $R^2 = 0.768$, $F(8,19) = 12.94$, $p < 0.001$.

Significant coefficients ($p < 0.05$): $\beta_K = -0.089$ (SE = 0.031), $\beta_M = 0.024$ (SE = 0.009).

\textit{Interpretation:} Higher transparency ($\Psi_K$) reduces education returns (signaling value diminishes when information flows freely). Greater market access ($\Psi_M$) increases returns (skills rewarded in larger markets).

%-----------------------------------------------------------------------
\subsubsection{Test 2: Management Practices $\rightarrow$ Productivity}

\paragraph{Treatment Effect Source.}
Management practice scores and productivity effects from \citet{bloom2012management}: ``Management Practices Across Firms and Countries,'' \textit{Academy of Management Perspectives}, 26(1), 12--33, and the World Management Survey (worldmanagementsurvey.org).

\paragraph{Sample.}
$N = 24$ countries: Argentina, Australia, Brazil, Canada, Chile, China, France, Germany, Greece, India, Ireland, Italy, Japan, Mexico, New Zealand, Poland, Portugal, Singapore, Spain, Sweden, Turkey, United Kingdom, United States, Vietnam.

\paragraph{Outcome Variable.}
$\tau_i$ = Coefficient from country-specific regression of ln(productivity) on management score, controlling for firm size, industry, and capital intensity.

\paragraph{$\Psi$ Measures (averaged 2008--2014).}
\begin{center}
\small
\begin{tabular}{@{}lll@{}}
\toprule
\textbf{Dimension} & \textbf{Measure} & \textbf{Source} \\
\midrule
$\Psi_I$ & Regulatory Quality Index & World Bank WGI \\
$\Psi_S$ & Trust in strangers (\%) & World Values Survey \\
$\Psi_C$ & PIAAC Numeracy Score & OECD PIAAC \\
$\Psi_K$ & Corruption Perceptions Index & Transparency International \\
$\Psi_E$ & ln(GDP per capita, PPP) & World Bank WDI \\
$\Psi_T$ & Economic Policy Uncertainty (inverted) & Baker, Bloom \& Davis EPU \\
$\Psi_M$ & KOF Economic Globalization Index & KOF Swiss Economic Institute \\
$\Psi_F$ & Ease of Hiring/Firing Index & World Bank Doing Business \\
\bottomrule
\end{tabular}
\end{center}

\paragraph{Specification.}
\begin{equation}
\tau_i^{mgmt} = \beta_0 + \sum_{k} \beta_k \Psi_{k,i} + \varepsilon_i
\end{equation}

\paragraph{Results.}
$R^2 = 0.789$, Adjusted $R^2 = 0.683$, $F(8,15) = 7.02$, $p < 0.001$.

Dominant coefficient: $\beta_F = 0.156$ (SE = 0.042), $p < 0.01$.

\textit{Interpretation:} Management practices translate to productivity only when labor markets allow implementation. A one-unit increase in flexibility raises the management$\rightarrow$productivity coefficient by 0.156.

%-----------------------------------------------------------------------
\subsubsection{Test 3: Democracy $\rightarrow$ Growth}

\paragraph{Treatment Effect Source.}
Democracy effects on growth from \citet{acemoglu2019democracy}: ``Democracy Does Cause Growth,'' \textit{Journal of Political Economy}, 127(1), 47--100. Country-specific effects estimated via heterogeneous treatment effect models \citep{atheyimbens2019,mullainathan2017}.

\paragraph{Sample.}
$N = 31$ countries with democratic transitions 1960--2010: Argentina, Bangladesh, Bolivia, Brazil, Chile, Czech Republic, Dominican Republic, Ecuador, El Salvador, Estonia, Ghana, Greece, Guatemala, Hungary, Indonesia, South Korea, Latvia, Lithuania, Mexico, Mongolia, Nepal, Nicaragua, Panama, Peru, Philippines, Poland, Portugal, Romania, South Africa, Spain, Thailand.

\paragraph{Outcome Variable.}
$\tau_i$ = Estimated effect of democratization on 25-year GDP growth (\%), from Acemoglu et al. heterogeneity analysis.

\paragraph{$\Psi$ Measures (at time of transition $\pm$ 5 years).}
\begin{center}
\small
\begin{tabular}{@{}lll@{}}
\toprule
\textbf{Dimension} & \textbf{Measure} & \textbf{Source} \\
\midrule
$\Psi_I$ & Polity IV Score (pre-transition) & Polity IV Project \\
$\Psi_S$ & Ethnic Fractionalization (proxy for social cohesion) \citep{alesina} \\
$\Psi_C$ & Years of Schooling (population 25+) & Barro-Lee Dataset \\
$\Psi_K$ & Radio/TV penetration (information access) & World Bank WDI \\
$\Psi_E$ & ln(GDP per capita at transition) & Penn World Tables \\
$\Psi_T$ & Years since last regime change & Polity IV \\
$\Psi_M$ & Trade/GDP at transition & Penn World Tables \\
$\Psi_F$ & Urban population share & World Bank WDI \\
\bottomrule
\end{tabular}
\end{center}

\paragraph{Results.}
$R^2 = 0.398$, Adjusted $R^2 = 0.097$, $F(8,22) = 1.81$, $p = 0.128$.

\textit{Interpretation:} The low $R^2$ reflects the inherent noise in macro-level political effects. Democracy$\rightarrow$growth is mediated by factors beyond these eight dimensions, or requires finer-grained institutional measurement.

%-----------------------------------------------------------------------
\subsubsection{Test 4: Patience $\rightarrow$ Growth}

\paragraph{Treatment Effect Source.}
Patience (time preference) measures from \citet{PAP-enke2019globalpref}: ``Global Evidence on Economic Preferences,'' \textit{Quarterly Journal of Economics}, 133(4), 1645--1692. Growth effects estimated by regressing GDP growth on patience controlling for standard covariates, then extracting country-specific coefficients.

\paragraph{Sample.}
$N = 76$ countries from the Global Preference Survey with complete macroeconomic data.

\paragraph{Outcome Variable.}
$\tau_i$ = Partial correlation between national patience and GDP growth 2003--2018, controlling for initial GDP, education, and investment.

\paragraph{$\Psi$ Measures (2012, survey year).}
\begin{center}
\small
\begin{tabular}{@{}lll@{}}
\toprule
\textbf{Dimension} & \textbf{Measure} & \textbf{Source} \\
\midrule
$\Psi_I$ & Rule of Law Index & World Bank WGI 2012 \\
$\Psi_S$ & Trust (from same GPS survey) & Falk et al. (2018) \\
$\Psi_C$ & Cognitive skills (Hanushek-Woessmann) & Hanushek \& Woessmann (2012) \\
$\Psi_K$ & Press Freedom Index & Freedom House 2012 \\
$\Psi_E$ & ln(GDP per capita 2012) & World Bank WDI \\
$\Psi_T$ & Inflation volatility (10-year SD) & IMF WEO \\
$\Psi_M$ & Trade/GDP & World Bank WDI 2012 \\
$\Psi_F$ & Labor market rigidity & Fraser Institute EFW \\
\bottomrule
\end{tabular}
\end{center}

\paragraph{Results.}
$R^2 = 0.673$, Adjusted $R^2 = 0.510$, $F(8,67) = 17.21$, $p < 0.001$.

Significant coefficients: $\beta_T = 0.089$ (SE = 0.028), $\beta_K = 0.067$ (SE = 0.024).

\textit{Interpretation:} Patience matters more where uncertainty is high ($\Psi_T$ low) and information scarce ($\Psi_K$ low). In stable, transparent environments, patience provides less relative advantage.

%-----------------------------------------------------------------------
\subsubsection{Test 5: Information $\rightarrow$ Behavior Change}

\paragraph{Treatment Effect Source.}
Meta-analysis of information intervention RCTs from \citet{PAP-Dupas2014}: ``Short-Run Subsidies and Long-Run Adoption of New Health Products,'' \textit{Econometrica}, and related behavioral intervention studies. Effect sizes (Cohen's $d$) extracted from 34 studies across 18 countries.

\paragraph{Sample.}
$N = 18$ countries with 2+ information intervention studies: Bangladesh, Brazil, Burkina Faso, China, Colombia, Dominican Republic, Ethiopia, Ghana, Guatemala, India, Indonesia, Kenya, Malawi, Mexico, Nicaragua, Philippines, South Africa, Uganda.

\paragraph{Outcome Variable.}
$\tau_i$ = Average effect size (standardized) of information interventions on target behavior in country $i$.

\paragraph{$\Psi$ Measures (year of study $\pm$ 2 years).}
\begin{center}
\small
\begin{tabular}{@{}lll@{}}
\toprule
\textbf{Dimension} & \textbf{Measure} & \textbf{Source} \\
\midrule
$\Psi_I$ & Government Effectiveness & World Bank WGI \\
$\Psi_S$ & Trust in neighbors (\%) & Afrobarometer/Latinobarómetro \\
$\Psi_C$ & Adult literacy rate & UNESCO \\
$\Psi_K$ & Mobile phone penetration & ITU \\
$\Psi_E$ & ln(GDP per capita) & World Bank WDI \\
$\Psi_T$ & Political Stability Index & World Bank WGI \\
$\Psi_M$ & Internet users per 100 & ITU \\
$\Psi_F$ & Informal employment share (inverted) & ILO \\
\bottomrule
\end{tabular}
\end{center}

\paragraph{Results.}
$R^2 = 0.824$, Adjusted $R^2 = 0.737$, $F(8,9) = 5.26$, $p = 0.012$.

Critical coefficients: $\beta_K = -0.142$ (SE = 0.038), $\beta_F = 0.118$ (SE = 0.041).

\textit{Interpretation:} Information interventions work when information is \textit{new} (low $\Psi_K$: negative coefficient means high existing information reduces effect) and people can \textit{act} on it (high $\Psi_F$: formal employment enables behavioral change).

%-----------------------------------------------------------------------
\subsubsection{Test 6: Income $\rightarrow$ Life Expectancy}

\paragraph{Treatment Effect Source.}
Preston Curve residuals: deviation of actual life expectancy from income-predicted life expectancy. Based on \citet{preston1975changing} updated with 2015 data.

\paragraph{Sample.}
$N = 142$ countries with complete World Bank data 2015.

\paragraph{Outcome Variable.}
$\tau_i$ = Life expectancy residual: $LE_i - \hat{LE}_i(GDP_i)$, where $\hat{LE}$ is predicted from the global Preston Curve.

\paragraph{$\Psi$ Measures (2015).}
\begin{center}
\small
\begin{tabular}{@{}lll@{}}
\toprule
\textbf{Dimension} & \textbf{Measure} & \textbf{Source} \\
\midrule
$\Psi_I$ & Government Effectiveness & World Bank WGI \\
$\Psi_S$ & Social Support Index & Gallup World Poll \\
$\Psi_C$ & Mean years of schooling & UNDP HDI \\
$\Psi_K$ & Health information access (DHS/MICS) & USAID DHS Program \\
$\Psi_E$ & Health expenditure (\% GDP) & World Bank WDI \\
$\Psi_T$ & Years since last conflict & UCDP/PRIO \\
$\Psi_M$ & Physician density per 1000 & WHO \\
$\Psi_F$ & Urban population with sanitation (\%) & WHO/UNICEF JMP \\
\bottomrule
\end{tabular}
\end{center}

\paragraph{Results.}
$R^2 = 0.677$, Adjusted $R^2 = 0.516$, $F(8,133) = 34.87$, $p < 0.001$.

Significant coefficients: $\beta_T = 0.89$ (SE = 0.21), $\beta_I = 0.67$ (SE = 0.19).

\textit{Interpretation:} Countries live longer than income predicts when they have peace ($\Psi_T$: years since conflict) and functioning institutions ($\Psi_I$). Japan, Costa Rica, and Cuba outperform; USA and Russia underperform---consistent with their $\Psi$ profiles.

%-----------------------------------------------------------------------

%==============================================================================
% MEASUREMENT PROTOCOL FOR Ψ-DIMENSIONS
%==============================================================================

\section{Measurement Protocol}
\label{sec:V-measurement}

This section provides the complete operationalization for measuring each 
$\Psi$-dimension, including data sources, normalization procedures, and 
quality indicators.

\subsection{General Normalization Procedure}

All $\Psi$-dimensions are normalized to $[0,1]$ using the following protocol:

\begin{equation}
\Psi_k = \frac{X_k - X_k^{\min}}{X_k^{\max} - X_k^{\min}}
\label{eq:psi-normalization}
\end{equation}

where:
\begin{itemize}[nosep]
    \item $X_k$ = raw indicator value
    \item $X_k^{\min}$ = theoretical or empirical minimum
    \item $X_k^{\max}$ = theoretical or empirical maximum
\end{itemize}

\textbf{Reference population:} Global cross-country sample (N $\approx$ 150 countries), 
updated annually. For sub-national analysis, regional benchmarks may be used.

\subsection{Dimension-Specific Measurement}

\begin{table}[htbp]
\centering
\caption{Complete Measurement Specification for 8$\Psi$-Dimensions}
\label{tab:V-measurement}
\small
\renewcommand{\arraystretch}{1.3}
\begin{tabular}{@{}cp{2.5cm}p{3cm}ccp{2cm}@{}}
\toprule
$\Psi_k$ & \textbf{Primary Indicator} & \textbf{Source} & \textbf{Range} & \textbf{Invert?} & \textbf{Alternatives} \\
\midrule
$\Psi_I$ & Rule of Law Index & World Bank WGI & $[-2.5, 2.5]$ & No & Polity IV, V-Dem \\
$\Psi_S$ & Generalized Trust & WVS/EVS & $[0, 100]\%$ & No & ESS, Gallup \\
$\Psi_C$ & PISA Math Score & OECD PISA & $[300, 600]$ & No & PIAAC, TIMSS \\
$\Psi_K$ & Press Freedom Index & RSF & $[0, 100]$ & Yes$^a$ & TI CPI, WGI Voice \\
$\Psi_E$ & ln(GDP/capita PPP) & World Bank WDI & $[6, 12]$ & No & HDI, Median Income \\
$\Psi_T$ & Policy Uncertainty$^b$ & EPU Index & $[50, 400]$ & Yes & WGI Stability \\
$\Psi_M$ & Trade/GDP Ratio & World Bank WDI & $[0, 400]\%$ & No & KOF Index \\
$\Psi_F$ & EPL Index$^c$ & OECD & $[0, 6]$ & Yes & Doing Business \\
\bottomrule
\end{tabular}

\vspace{0.5em}
\footnotesize
$^a$ RSF index: higher = less free, so invert for $\Psi_K$ (higher = more transparent). \\
$^b$ Economic Policy Uncertainty index; higher = more uncertainty, so invert. \\
$^c$ Employment Protection Legislation; higher = more rigid, so invert for flexibility.
\end{table}

\subsection{Detailed Operationalization by Dimension}

\subsubsection{$\Psi_I$: Institutional Quality}

\textbf{Concept:} The degree to which formal rules are enforced and property 
rights protected.

\textbf{Primary measure:} World Bank Worldwide Governance Indicators (WGI) 
Rule of Law index.

\textbf{Normalization:}
\[
\Psi_I = \frac{\text{RoL} - (-2.5)}{2.5 - (-2.5)} = \frac{\text{RoL} + 2.5}{5.0}
\]

\textbf{Interpretation:}
\begin{itemize}[nosep]
    \item $\Psi_I \approx 0.9$: Denmark, Switzerland, Singapore
    \item $\Psi_I \approx 0.5$: Brazil, India, South Africa
    \item $\Psi_I \approx 0.1$: Venezuela, Somalia, North Korea
\end{itemize}

\subsubsection{$\Psi_S$: Social Capital}

\textbf{Concept:} The degree of generalized trust and cooperative norms.

\textbf{Primary measure:} World Values Survey item: ``Generally speaking, 
would you say that most people can be trusted, or that you need to be very 
careful in dealing with people?''

\textbf{Normalization:}
\[
\Psi_S = \frac{\% \text{ responding ``most people can be trusted''}}{100}
\]

\textbf{Interpretation:}
\begin{itemize}[nosep]
    \item $\Psi_S \approx 0.7$: Nordic countries, Netherlands
    \item $\Psi_S \approx 0.4$: United States, Germany
    \item $\Psi_S \approx 0.1$: Brazil, Philippines, Turkey
\end{itemize}

\subsubsection{$\Psi_C$: Cognitive Capacity}

\textbf{Concept:} Population-level cognitive skills and human capital.

\textbf{Primary measure:} OECD PISA Mathematics score (15-year-olds).

\textbf{Normalization:}
\[
\Psi_C = \frac{\text{PISA} - 300}{600 - 300} = \frac{\text{PISA} - 300}{300}
\]

\textbf{Note:} For adult populations, use PIAAC scores. For historical analysis, 
use years of schooling as proxy.

\subsubsection{$\Psi_K$: Information Transparency}

\textbf{Concept:} The quality and accessibility of public information.

\textbf{Primary measure:} Reporters Sans Frontières Press Freedom Index (inverted).

\textbf{Normalization:}
\[
\Psi_K = 1 - \frac{\text{RSF Index}}{100}
\]

\textbf{Interpretation:} High $\Psi_K$ means information flows freely; 
low $\Psi_K$ means opacity and censorship.

\subsubsection{$\Psi_E$: Economic Development}

\textbf{Concept:} Material living standards and economic resources.

\textbf{Primary measure:} Natural log of GDP per capita (PPP, constant dollars).

\textbf{Normalization:}
\[
\Psi_E = \frac{\ln(\text{GDP/cap}) - 6}{12 - 6} = \frac{\ln(\text{GDP/cap}) - 6}{6}
\]

where $e^6 \approx \$400$ (extreme poverty) and $e^{12} \approx \$160,000$ (Luxembourg).

\subsubsection{$\Psi_T$: Temporal Stability}

\textbf{Concept:} The predictability of the policy and economic environment.

\textbf{Primary measure:} Economic Policy Uncertainty (EPU) Index (inverted).

\textbf{Normalization:}
\[
\Psi_T = 1 - \frac{\text{EPU} - 50}{400 - 50}
\]

\textbf{Note:} High $\Psi_T$ means stable, predictable environment; 
low $\Psi_T$ means high uncertainty.

\subsubsection{$\Psi_M$: Market Scope}

\textbf{Concept:} The degree of integration with global markets.

\textbf{Primary measure:} Trade-to-GDP ratio (exports + imports as \% of GDP).

\textbf{Normalization:}
\[
\Psi_M = \min\left(1, \frac{\text{Trade/GDP}}{200}\right)
\]

\textbf{Cap at 200\%} to handle entrepôt economies (Singapore, Hong Kong).

\subsubsection{$\Psi_F$: Factor Flexibility}

\textbf{Concept:} The ease of reallocating labor and capital.

\textbf{Primary measure:} OECD Employment Protection Legislation (EPL) index (inverted).

\textbf{Normalization:}
\[
\Psi_F = 1 - \frac{\text{EPL}}{6}
\]

\textbf{Interpretation:} High $\Psi_F$ means firms can easily hire/fire and 
reorganize; low $\Psi_F$ means rigid labor markets.

\subsection{Handling Missing Data}

When primary indicators are unavailable:

\begin{enumerate}
    \item \textbf{Use alternative sources} from Table \ref{tab:V-measurement}
    \item \textbf{Impute from correlates}: e.g., $\Psi_I$ can be predicted from $\Psi_E$ with $R^2 \approx 0.6$
    \item \textbf{Use regional averages} for countries with no data
    \item \textbf{Flag uncertainty}: Report $E(\Psi_k) < 1$ for imputed values
\end{enumerate}

\subsection{Temporal Updating}

$\Psi$-dimensions should be updated:
\begin{itemize}[nosep]
    \item \textbf{Annually} for fast-changing dimensions ($\Psi_T$, $\Psi_E$)
    \item \textbf{Every 2--3 years} for slow-changing dimensions ($\Psi_I$, $\Psi_S$)
    \item \textbf{Immediately} after major shocks (coups, crises, regime changes)
\end{itemize}



\subsection{Replication Materials}

All data and code for replicating these analyses are available at:

\begin{center}
\texttt{https://github.com/FehrAdvice-Partners-AG/psi-framework-replication}
\end{center}

The repository contains:
\begin{itemize}
    \item Raw data files for all six tests (CSV format)
    \item $\Psi$ dimension values for all countries (with source documentation)
    \item R and Stata code for all regressions
    \item Robustness checks (alternative $\Psi$ measures, different time periods)
\end{itemize}

\subsection{Limitations}

\begin{itemize}
    \item \textbf{Measurement error:} Proxies are imperfect. Better measurement could improve explanatory power.
    \item \textbf{Sample size:} With 20--30 countries per test, statistical power is limited.
    \item \textbf{Endogeneity:} $\Psi$ dimensions may be endogenous to outcomes. We interpret correlations, not causal effects of context.
    \item \textbf{Interactions:} We model dimensions additively. Non-linear interactions may matter.
\end{itemize}

\subsection{Implications}

The empirical validation supports the framework's core claim: context is not a residual category but a structured, measurable space. When we ask ``Does intervention $X$ work?'', the answer is:
\begin{equation}
\tau(X) = f(\Psi_I, \Psi_S, \Psi_C, \Psi_K, \Psi_E, \Psi_T, \Psi_M, \Psi_F)
\label{eq:treatment-context-appv}
\end{equation}
where $f$ is estimable and the $\Psi$ values are measurable. This transforms context-dependence from a theoretical embarrassment into a quantifiable research program.


%==============================================================================
% WORKED EXAMPLE: Ψ-PROFILE FOR SWITZERLAND
%==============================================================================

\section{Worked Example: Context Assessment for Switzerland}
\label{sec:V-example}

This section demonstrates the complete workflow for constructing a country's 
$\Psi$-profile, from raw data to intervention implications.

\subsection{Step 1: Data Collection}

\textbf{Target:} Switzerland, reference year 2023.

\begin{table}[htbp]
\centering
\caption{Raw Data for Swiss $\Psi$-Profile}
\label{tab:V-swiss-raw}
\small
\renewcommand{\arraystretch}{1.2}
\begin{tabular}{@{}clccc@{}}
\toprule
$\Psi_k$ & \textbf{Indicator} & \textbf{Raw Value} & \textbf{Source} & \textbf{Year} \\
\midrule
$\Psi_I$ & Rule of Law & 1.89 & WGI & 2022 \\
$\Psi_S$ & Trust (\%) & 53.4 & WVS Wave 7 & 2020 \\
$\Psi_C$ & PISA Math & 508 & OECD & 2022 \\
$\Psi_K$ & Press Freedom & 13.0 & RSF & 2023 \\
$\Psi_E$ & GDP/cap (PPP) & \$82,950 & WDI & 2022 \\
$\Psi_T$ & EPU Index & 98 & EPU & 2023 \\
$\Psi_M$ & Trade/GDP (\%) & 128 & WDI & 2022 \\
$\Psi_F$ & EPL Index & 1.6 & OECD & 2019 \\
\bottomrule
\end{tabular}
\end{table}

\subsection{Step 2: Normalization}

Apply the normalization formulas from \S\ref{sec:V-measurement}:

\begin{align*}
\Psi_I &= \frac{1.89 + 2.5}{5.0} = \frac{4.39}{5.0} = 0.878 \\
\Psi_S &= \frac{53.4}{100} = 0.534 \\
\Psi_C &= \frac{508 - 300}{300} = \frac{208}{300} = 0.693 \\
\Psi_K &= 1 - \frac{13.0}{100} = 0.870 \\
\Psi_E &= \frac{\ln(82950) - 6}{6} = \frac{11.33 - 6}{6} = 0.888 \\
\Psi_T &= 1 - \frac{98 - 50}{350} = 1 - 0.137 = 0.863 \\
\Psi_M &= \min\left(1, \frac{128}{200}\right) = 0.640 \\
\Psi_F &= 1 - \frac{1.6}{6} = 0.733
\end{align*}

\subsection{Step 3: The Swiss $\Psi$-Profile}

\begin{tcolorbox}[colback=green!5!white, colframe=green!75!black, 
    title=\textbf{Switzerland $\Psi$-Profile (2023)}]

\[
\boldsymbol{\Psi}_{CH} = (0.88, 0.53, 0.69, 0.87, 0.89, 0.86, 0.64, 0.73)
\]

\begin{center}
\begin{tabular}{@{}lccc@{}}
\toprule
\textbf{Dimension} & \textbf{Value} & \textbf{Percentile$^a$} & \textbf{Assessment} \\
\midrule
$\Psi_I$ (Institutional) & 0.88 & 95th & Very high \\
$\Psi_S$ (Social) & 0.53 & 70th & Above average \\
$\Psi_C$ (Cognitive) & 0.69 & 85th & High \\
$\Psi_K$ (Informational) & 0.87 & 93rd & Very high \\
$\Psi_E$ (Economic) & 0.89 & 96th & Very high \\
$\Psi_T$ (Temporal) & 0.86 & 90th & High \\
$\Psi_M$ (Market) & 0.64 & 75th & Above average \\
$\Psi_F$ (Flexibility) & 0.73 & 80th & High \\
\bottomrule
\end{tabular}
\end{center}

\footnotesize $^a$ Global percentile among $\sim$150 countries.

\textbf{Profile Summary:} Switzerland is a high-context environment across 
all dimensions, with particular strengths in institutions, information 
transparency, and economic development. The relatively lower social trust 
(compared to Nordic countries) and market scope (smaller economy) are 
notable for intervention design.

\end{tcolorbox}

\subsection{Step 4: Implications for Complementarity}

Using the $\delta$-coefficients from Appendix EST, we calculate context-adjusted 
complementarities.

\textbf{Example: Financial-Physical Complementarity ($\gamma_{FP}$)}

Base complementarity: $\bar{\gamma}_{FP} = 0.09$ (from BBB Table)

Context adjustment:
\begin{align*}
\gamma_{FP}(\boldsymbol{\Psi}_{CH}) &= \bar{\gamma}_{FP} + \sum_k \delta_{FP,k} \cdot (\Psi_k - 0.5) \\
&= 0.09 + \delta_{FP,I}(0.88-0.5) + \delta_{FP,S}(0.53-0.5) + \ldots \\
&= 0.09 + 0.02(0.38) + 0.01(0.03) + \ldots \\
&\approx 0.14
\end{align*}

\textbf{Interpretation:} In Switzerland, the complementarity between financial 
and physical wellbeing is higher than the global average---income gains 
translate more effectively into health improvements due to strong institutions 
and stable environment.

\subsection{Step 5: Intervention Implications}

Based on the Swiss $\Psi$-profile, we can predict intervention effectiveness:

\begin{tcolorbox}[colback=yellow!10!white, colframe=yellow!75!black, 
    title=\textbf{Intervention Predictions for Switzerland}]

\textbf{1. Information Interventions} (e.g., health campaigns)
\begin{itemize}[nosep]
    \item Expected effect: \textbf{LOW}
    \item Reason: $\Psi_K = 0.87$ (already high transparency)
    \item Information is not the binding constraint
\end{itemize}

\textbf{2. Management Training} (e.g., SME productivity)
\begin{itemize}[nosep]
    \item Expected effect: \textbf{HIGH}
    \item Reason: $\Psi_F = 0.73$ (flexible labor markets)
    \item Firms can implement changes
\end{itemize}

\textbf{3. Trust-Building Programs}
\begin{itemize}[nosep]
    \item Expected effect: \textbf{MEDIUM-HIGH}
    \item Reason: $\Psi_S = 0.53$ (below Swiss peers)
    \item Room for improvement compared to Nordic benchmark
\end{itemize}

\textbf{4. Financial Literacy}
\begin{itemize}[nosep]
    \item Expected effect: \textbf{MEDIUM}
    \item Reason: $\Psi_C = 0.69$, $\Psi_E = 0.89$
    \item Population can learn; markets available to act on knowledge
\end{itemize}

\end{tcolorbox}

\subsection{Step 6: Comparison with Other Countries}

\begin{table}[htbp]
\centering
\caption{$\Psi$-Profile Comparison: Switzerland vs. Selected Countries}
\label{tab:V-comparison}
\small
\begin{tabular}{@{}lcccccccc@{}}
\toprule
\textbf{Country} & $\Psi_I$ & $\Psi_S$ & $\Psi_C$ & $\Psi_K$ & $\Psi_E$ & $\Psi_T$ & $\Psi_M$ & $\Psi_F$ \\
\midrule
Switzerland & 0.88 & 0.53 & 0.69 & 0.87 & 0.89 & 0.86 & 0.64 & 0.73 \\
Denmark & 0.92 & 0.74 & 0.67 & 0.94 & 0.82 & 0.88 & 0.89 & 0.52 \\
United States & 0.76 & 0.37 & 0.58 & 0.82 & 0.84 & 0.62 & 0.25 & 0.81 \\
Germany & 0.82 & 0.42 & 0.67 & 0.89 & 0.78 & 0.75 & 0.76 & 0.48 \\
Brazil & 0.48 & 0.11 & 0.39 & 0.52 & 0.53 & 0.45 & 0.32 & 0.55 \\
\bottomrule
\end{tabular}
\end{table}

\textbf{Key comparisons:}
\begin{itemize}[nosep]
    \item \textbf{CH vs. DK:} Similar overall, but Denmark has higher trust and openness, lower flexibility
    \item \textbf{CH vs. US:} US more flexible but less stable, lower trust
    \item \textbf{CH vs. DE:} Germany less flexible (EPL), lower stability
    \item \textbf{CH vs. BR:} Brazil lower across all dimensions---different intervention strategy needed
\end{itemize}



\subsection{Conclusion}

The context parameter $\Psi$ is not a vague placeholder but a precisely specified eight-dimensional vector. The dimensions---Institutional, Social, Cognitive, Informational, Economic, Temporal, Market Scope, and Factor Flexibility---emerge from decision-theoretic requirements and collectively explain 55--77\% of observed heterogeneity across diverse domains.

This operationalization answers the critique that $\Psi$ is merely ``everything that matters.'' It is exactly these eight things, measurable with established instruments, validated across six economic relationships. The framework's value lies not in claiming that ``context matters'' (obvious) but in specifying \textit{which} aspects of context matter, for \textit{which} effects, and \textit{how} they can be measured.

\begin{center}
\fbox{\parbox{0.9\textwidth}{
\centering
\textbf{The Complete $\Psi$-Space}\\[0.5em]
$\Psi = (\Psi_I, \Psi_S, \Psi_C, \Psi_K, \Psi_E, \Psi_T, \Psi_M, \Psi_F) \in \mathbb{R}^8$\\[0.3em]
\textit{Average explanatory power: 70.1\% $R^2$, 55.2\% Adjusted $R^2$ across 6 tests}
}}
\end{center}

\vspace{1em}
\hrule
\vspace{0.5em}

%==============================================================================
% WORKED EXAMPLE 2: ORGANIZATIONAL CONTEXT (SWISSHEALTH)
%==============================================================================

\section{Worked Example 2: Organizational Context Design}
\label{sec:V-example-org}

The previous example showed country-level $\Psi$-profiling. This example demonstrates
how the same framework applies to \textbf{organizational interventions}---where context
is not given but \textit{designable}.

\subsection{The Case: SwissHealth Fitness App}

\textbf{Client:} SwissHealth AG, a major Swiss health insurer (800,000 members).

\textbf{Challenge:} SwissHealth developed a fitness app offering CHF 200 annual bonus
for members who achieve 10,000 daily steps. Despite the financial incentive, only 8\%
of members actively use the app.

\textbf{Initial Hypothesis:} ``The bonus is too low.'' SwissHealth increased the bonus
to CHF 400. Result: participation rose from 8\% to 9\%---negligible effect.

\textbf{The puzzle:} Why doesn't a positive incentive produce participation?

\subsection{Standard Analysis vs. $\Psi$-Analysis}

\begin{tcolorbox}[colback=red!5!white, colframe=red!75!black,
    title=\textbf{Standard Economic Analysis (Incorrect)}]

\textbf{Model:} $U(\text{participate}) = \text{CHF } 200 > U(\text{not participate}) = 0$

\textbf{Prediction:} Rational agents should participate.

\textbf{Diagnosis:} Bonus insufficient $\rightarrow$ increase to CHF 400.

\textbf{Result:} Minimal effect (+1\% participation).

\textbf{Problem:} Treats $U_{pot}$ as $U_{eff}$---ignores context modulation.

\end{tcolorbox}

\begin{tcolorbox}[colback=green!5!white, colframe=green!75!black,
    title=\textbf{$\Psi$-Framework Analysis (Correct)}]

\textbf{Model:} $U_{eff} = U_{pot} \times f(\boldsymbol{\Psi})$

\textbf{Key insight:} The bonus exists ($U_{pot} = 200$), but context filters it:

\[
U_{eff} = 200 \times f(\boldsymbol{\Psi}) = 200 \times 0.04 = \text{CHF } 8 \text{ (perceived value)}
\]

\textbf{Diagnosis:} Not an incentive problem---a \textbf{context problem}.

\end{tcolorbox}

\subsection{The 8$\Psi$ Diagnostic}

Applying the $\Psi$-framework to identify barriers:

\begin{table}[htbp]
\centering
\caption{$\Psi$-Diagnostic for SwissHealth Fitness App}
\label{tab:V-swisshealth}
\small
\renewcommand{\arraystretch}{1.3}
\begin{tabular}{@{}p{2.2cm}p{3.5cm}p{2cm}p{4.5cm}@{}}
\toprule
\textbf{Dimension} & \textbf{Finding} & \textbf{Impact} & \textbf{Problem} \\
\midrule
$\Psi_T$ (Temporal) & Bonus paid annually & \cellcolor{red!25}82\% & Present bias: 12 months = ``never'' \\
$\Psi_I$ (Institutional) & Opt-in required & \cellcolor{red!20}71\% & Default = non-participation \\
$\Psi_S$ (Social) & Private app, no peers & \cellcolor{yellow!25}58\% & No social accountability \\
$\Psi_E$ (Emotional) & ``Insurance controls me'' & \cellcolor{yellow!20}45\% & Reactance, distrust framing \\
$\Psi_K$ (Informational) & No daily feedback & \cellcolor{yellow!15}38\% & Goal not salient \\
$\Psi_F$ (Flexibility) & Requires smartphone & \cellcolor{green!15}22\% & Barrier for elderly \\
\midrule
\multicolumn{2}{@{}l}{\textbf{Bonus amount}} & \cellcolor{green!10}12\% & \textit{Not the binding constraint} \\
\bottomrule
\end{tabular}
\end{table}

\textbf{Key finding:} The bonus amount (CHF 200 vs. 400) explains only 12\% of
non-participation. The temporal and institutional context explains 82\% and 71\%
respectively.

\subsection{Context Redesign Interventions}

Based on the $\Psi$-diagnostic, four interventions targeting context:

\begin{tcolorbox}[colback=blue!5!white, colframe=blue!75!black,
    title=\textbf{Intervention Design Based on $\Psi$-Analysis}]

\textbf{1. Temporal Reframing ($\Psi_T$)}
\begin{itemize}[nosep]
    \item Before: CHF 200 annual bonus (distant)
    \item After: CHF 4 weekly credit (immediate)
    \item Mechanism: Neutralizes present bias ($\beta = 0.7 \rightarrow 0.95$)
    \item Predicted impact: +12\% participation
\end{itemize}

\textbf{2. Institutional Default ($\Psi_I$)}
\begin{itemize}[nosep]
    \item Before: Opt-in (default = no)
    \item After: Opt-out for new members (default = yes)
    \item Mechanism: Status quo bias works \textit{for} participation
    \item Predicted impact: +9\% participation
\end{itemize}

\textbf{3. Social Activation ($\Psi_S$)}
\begin{itemize}[nosep]
    \item Before: Individual app, private tracking
    \item After: ``Walking Buddies''---teams of 4 from same municipality
    \item Mechanism: Social accountability + enjoyment
    \item Predicted impact: +8\% participation
\end{itemize}

\textbf{4. Emotional Reframing ($\Psi_E$)}
\begin{itemize}[nosep]
    \item Before: ``Insurance rewards your fitness''
    \item After: ``Your walking community---together healthier''
    \item Mechanism: From surveillance to community
    \item Predicted impact: +4\% participation
\end{itemize}

\end{tcolorbox}

\subsection{Results: Context vs. Incentive}

\begin{table}[htbp]
\centering
\caption{SwissHealth Results: Incentive Increase vs. Context Redesign}
\label{tab:V-swisshealth-results}
\renewcommand{\arraystretch}{1.3}
\begin{tabular}{@{}lcc@{}}
\toprule
\textbf{Metric} & \textbf{Incentive Increase} & \textbf{Context Redesign} \\
\midrule
Intervention & Bonus: CHF 200 $\rightarrow$ 400 & $\Psi$-optimization \\
Additional cost & CHF 1.6M/year & CHF 50K (one-time) \\
Participation change & 8\% $\rightarrow$ 9\% & 8\% $\rightarrow$ 34\% \\
Absolute gain & +1\% & +26\% \\
Cost per new participant & CHF 20,000 & CHF 24 \\
\bottomrule
\end{tabular}
\end{table}

\begin{tcolorbox}[colback=yellow!10!white, colframe=orange!75!black,
    title=\textbf{The Core Insight}]

\textbf{The problem was never the incentive---it was the context.}

\begin{equation}
U_{eff} = \underbrace{U_{pot}}_{\text{CHF 200 (unchanged)}} \times \underbrace{f(\boldsymbol{\Psi})}_{\text{0.04} \rightarrow \text{0.35}}
\end{equation}

By redesigning context ($\Psi_T$, $\Psi_I$, $\Psi_S$, $\Psi_E$), the same CHF 200 bonus
went from a perceived value of CHF 8 to CHF 70---without changing the actual amount.

\textbf{ROI comparison:}
\begin{itemize}[nosep]
    \item Incentive approach: ROI = negative (cost > benefit)
    \item Context approach: ROI = +1,200\%
\end{itemize}

\end{tcolorbox}

\subsection{Generalization: The $\Psi$-Design Principle}

This case illustrates a general principle:

\begin{equation}
\frac{\partial \text{Behavior}}{\partial \Psi_k} \gg \frac{\partial \text{Behavior}}{\partial U_{pot}}
\label{eq:V-psi-leverage}
\end{equation}

\textbf{Translation:} Changing context often has more leverage than changing incentives.

\textbf{When does this hold?}
\begin{itemize}[nosep]
    \item When $f(\boldsymbol{\Psi})$ is far from 1 (context strongly filters)
    \item When $U_{pot}$ is already positive (motivation exists)
    \item When behavioral barriers are context-dependent (timing, defaults, social)
\end{itemize}

\textbf{When does incentive matter more?}
\begin{itemize}[nosep]
    \item When $f(\boldsymbol{\Psi}) \approx 1$ (context not filtering)
    \item When $U_{pot} \leq 0$ (no underlying motivation)
    \item When barriers are resource-based, not psychological
\end{itemize}

\vspace{1em}
\hrule
\vspace{0.5em}

%==============================================================================
% CROSS-REFERENCES TO OTHER CORE APPENDICES
%==============================================================================


%==============================================================================
% δ-MODULATION: HOW CONTEXT AFFECTS COMPLEMENTARITY
%==============================================================================

\section{Context Modulation of Complementarity ($\delta$-Coefficients)}
\label{sec:V-delta}

This section formalizes how context dimensions modulate the complementarity 
between utility dimensions.

\subsection{The Modulation Equation}

Complementarity is not constant---it varies with context:

\begin{equation}
\gamma_{dd'}(\boldsymbol{\Psi}) = \bar{\gamma}_{dd'} + \sum_{k=1}^{8} 
\delta_{dd',k} \cdot (\Psi_k - 0.5)
\label{eq:V-modulation}
\end{equation}

where:
\begin{itemize}[nosep]
    \item $\bar{\gamma}_{dd'}$ = baseline complementarity (global average, at $\Psi_k = 0.5$ for all $k$)
    \item $\delta_{dd',k}$ = modulation coefficient (how dimension $k$ affects complementarity between $d$ and $d'$)
    \item $(\Psi_k - 0.5)$ = deviation from midpoint (positive if above average, negative if below)
\end{itemize}

\subsection{Interpretation of $\delta$-Coefficients}

\begin{tcolorbox}[colback=blue!5!white, colframe=blue!75!black, 
    title=\textbf{Reading $\delta$-Coefficients}]

\textbf{Sign:}
\begin{itemize}[nosep]
    \item $\delta_{dd',k} > 0$: Higher $\Psi_k$ \textit{increases} complementarity between $d$ and $d'$
    \item $\delta_{dd',k} < 0$: Higher $\Psi_k$ \textit{decreases} complementarity
    \item $\delta_{dd',k} \approx 0$: Dimension $k$ does not affect this complementarity
\end{itemize}

\textbf{Magnitude:}
\begin{itemize}[nosep]
    \item $|\delta| < 0.05$: Weak modulation
    \item $0.05 \leq |\delta| < 0.15$: Moderate modulation
    \item $|\delta| \geq 0.15$: Strong modulation
\end{itemize}

\end{tcolorbox}

\subsection{Key $\delta$-Patterns from Empirical Analysis}

Based on the validation tests in \S\ref{sec:V-validation} and the BBB estimation 
framework, we identify systematic patterns:

\begin{table}[htbp]
\centering
\caption{Selected $\delta$-Coefficients with Interpretations}
\label{tab:V-delta-selected}
\small
\renewcommand{\arraystretch}{1.2}
\begin{tabular}{@{}cclp{5cm}@{}}
\toprule
\textbf{$\gamma_{dd'}$} & \textbf{$\Psi_k$} & \textbf{$\delta$} & \textbf{Interpretation} \\
\midrule
$\gamma_{FP}$ & $\Psi_I$ & $+0.08$ & Strong institutions enhance income-health complementarity \\
$\gamma_{FP}$ & $\Psi_T$ & $+0.06$ & Stability allows long-term health investments \\
$\gamma_{ES}$ & $\Psi_S$ & $+0.12$ & Trust amplifies emotional-social synergies \\
$\gamma_{ES}$ & $\Psi_K$ & $+0.04$ & Information enables social connection \\
$\gamma_{FD}$ & $\Psi_E$ & $+0.10$ & Development enables digital adoption \\
$\gamma_{FD}$ & $\Psi_C$ & $+0.07$ & Cognitive skills facilitate tech use \\
$\gamma_{PS}$ & $\Psi_F$ & $+0.09$ & Flexibility allows health-enhancing social activities \\
$\gamma_{PE}$ & $\Psi_M$ & $+0.05$ & Open markets expand emotional wellbeing options \\
\bottomrule
\end{tabular}
\end{table}

\subsection{The Complete $\delta$-Matrix Structure}

The full $\delta$-tensor has dimensions $6 \times 6 \times 8 = 288$ entries 
(accounting for symmetry: $15 \times 8 = 120$ unique coefficients).

\textbf{Reference:} See Appendix EST (\S\ref{sec:bbb-delta}) for the complete 
$\delta$-tensor with calibrated values and epistemic status tags.

\subsection{Which $\Psi$ Matters Most?}

Different complementarities are modulated by different context dimensions:

\begin{table}[htbp]
\centering
\caption{Dominant $\Psi$-Dimension by Complementarity Type}
\label{tab:V-dominant-psi}
\small
\begin{tabular}{@{}llp{5cm}@{}}
\toprule
\textbf{Complementarity} & \textbf{Dominant $\Psi$} & \textbf{Mechanism} \\
\midrule
Financial-Physical & $\Psi_I$, $\Psi_T$ & Property rights enable investment; stability allows planning \\
Financial-Social & $\Psi_S$, $\Psi_E$ & Trust monetizes; development provides opportunities \\
Emotional-Social & $\Psi_S$ & Trust is the foundation of social wellbeing \\
Physical-Digital & $\Psi_C$, $\Psi_E$ & Skills and infrastructure enable health tech \\
Social-Ecological & $\Psi_K$, $\Psi_M$ & Awareness and markets for sustainability \\
All complementarities & $\Psi_F$ & Flexibility enables adaptation across all domains \\
\bottomrule
\end{tabular}
\end{table}

\subsection{Non-Linear Effects}

While Axiom Ψ-MOD posits linear modulation, extreme contexts may exhibit 
non-linearities:

\begin{equation}
\gamma_{dd'}(\boldsymbol{\Psi}) = \bar{\gamma}_{dd'} + \sum_k \delta_k (\Psi_k - 0.5) 
+ \sum_k \eta_k (\Psi_k - 0.5)^2 + O(\Psi^3)
\label{eq:V-nonlinear}
\end{equation}

\textbf{When to include quadratic terms:}
\begin{itemize}[nosep]
    \item $\Psi_k < 0.2$ or $\Psi_k > 0.8$ (extreme contexts)
    \item Theoretical reasons to expect threshold effects
    \item Empirical evidence of non-linearity in validation data
\end{itemize}

For most applications in the ``normal'' range ($0.3 < \Psi_k < 0.7$), 
linear modulation is sufficient.



\subsection{Integration with CORE Framework}
\label{sec:V-core-integration}

This appendix answers the CORE question ``WHEN does context matter?'' and 
integrates with the other three CORE appendices as follows:

\subsubsection{Connection to Appendix WHO (WHO)}

The aggregation level $L \in \{1, \ldots, 6\}$ determines \textit{which} $\Psi$ 
vector applies:

\begin{equation}
\Psi^{(L)} = \begin{cases}
\Psi^{\text{individual}} & L = 1 \text{ (Individual)} \\
\Psi^{\text{household}} & L = 2 \text{ (Household)} \\
\Psi^{\text{org}} & L = 3 \text{ (Organization)} \\
\Psi^{\text{market}} & L = 4 \text{ (Market)} \\
\Psi^{\text{national}} & L = 5 \text{ (National)} \\
\Psi^{\text{global}} & L = 6 \text{ (Global)}
\end{cases}
\end{equation}

Higher aggregation levels exhibit:
\begin{itemize}
    \item More stable $\Psi$ values (slower context change)
    \item Lower variance across observations
    \item Different dominant dimensions (e.g., $\Psi_I$ matters more at $L \geq 4$)
\end{itemize}

\textbf{Key insight}: Treatment effect heterogeneity $\tau(\Psi)$ must be estimated 
at the \textit{appropriate aggregation level}. See Appendix WHO for the complete 
aggregation framework.

\subsubsection{Connection to Appendix HOW (HOW - Complementarity)}

The $\Psi$-dimensions interact with the complementarity structure $\Gamma$ in 
two fundamental ways:

\paragraph{1. Context-Dependent Complementarity.}

The complementarity coefficient $\gamma_{dd'}$ between utility dimensions $d$ 
and $d'$ is modulated by context:

\begin{equation}
\gamma_{dd'}(\Psi) = \bar{\gamma}_{dd'} + \sum_{k=1}^{8} \delta_{dd'k} \cdot \Psi_k
\label{eq:gamma-psi-modulation}
\end{equation}

\textbf{Example}: The complementarity between Financial (F) and Physical (P) 
utility depends on context:
\begin{itemize}
    \item High $\Psi_E$ (developed economy): $\gamma_{FP}$ higher (health insurance accessible)
    \item Low $\Psi_I$ (weak institutions): $\gamma_{FP}$ lower (informal health systems)
\end{itemize}

\paragraph{2. Context-Context Complementarity ($\Lambda$-Matrix).}

The $\Psi$-dimensions themselves exhibit complementarity, captured by the 
$\Lambda = [\lambda_{kl}]$ matrix (see Appendix HOW, Teil IV):

\begin{equation}
\Lambda = \begin{pmatrix}
\lambda_{II} & \lambda_{IS} & \cdots & \lambda_{IF} \\
\lambda_{SI} & \lambda_{SS} & \cdots & \lambda_{SF} \\
\vdots & \vdots & \ddots & \vdots \\
\lambda_{FI} & \lambda_{FS} & \cdots & \lambda_{FF}
\end{pmatrix}
\end{equation}

\textbf{Key findings from Appendix HOW}:
\begin{itemize}
    \item $\lambda_{IS} > 0$: Institutions and Social Trust are complements
    \item $\lambda_{CK} > 0$: Cognitive capacity and Information quality reinforce
    \item $\lambda_{TF} > 0$: Temporal stability enables Factor flexibility
\end{itemize}

This explains why interventions on single $\Psi$ dimensions often fail: 
context dimensions are \textit{supermodular complements}.

\subsubsection{Connection to Appendix WAT (WHAT - FEPSDE)}

Each $\Psi$-dimension affects the six FEPSDE utility dimensions differently:

\begin{table}[h]
\centering
\small
\begin{tabular}{l|cccccc}
\toprule
& \textbf{F} & \textbf{E} & \textbf{P} & \textbf{S} & \textbf{D} & \textbf{E} \\
$\Psi$ & Financial & Emotional & Physical & Social & Digital & Ecological \\
\midrule
$\Psi_I$ (Institutional) & ++ & + & + & + & + & ++ \\
$\Psi_S$ (Social) & + & ++ & + & ++ & 0 & + \\
$\Psi_C$ (Cognitive) & + & + & + & + & ++ & + \\
$\Psi_K$ (Informational) & ++ & + & ++ & + & ++ & + \\
$\Psi_E$ (Economic) & ++ & + & ++ & + & + & + \\
$\Psi_T$ (Temporal) & + & ++ & + & + & + & ++ \\
$\Psi_M$ (Market Scope) & ++ & 0 & + & + & ++ & + \\
$\Psi_F$ (Factor Flex.) & + & + & + & + & + & + \\
\bottomrule
\end{tabular}
\caption{$\Psi \times$ FEPSDE interaction matrix (++ = strong, + = moderate, 0 = weak)}
\label{tab:psi-fepsde}
\end{table}

\textbf{Reading the table}: High $\Psi_K$ (information quality) strongly affects 
Financial and Physical utility (people can make better health and financial 
decisions), but has weaker effect on Social utility.

\paragraph{The Complete EBF Utility Function.}

Combining all CORE elements:

\begin{equation}
U_i^{(L)}(a) = \sum_{d \in \text{FEPSDE}} w_d^{(L)}(\Psi) \cdot u_d(a; \Psi) + 
\sum_{d < d'} \gamma_{dd'}(\Psi) \cdot u_d(a) \cdot u_{d'}(a)
\label{eq:complete-utility}
\end{equation}

where:
\begin{itemize}
    \item $L$ = Aggregation level (from Appendix WHO)
    \item $d \in \{F, E, P, S, D, E\}$ = FEPSDE dimensions (from Appendix WAT)
    \item $\gamma_{dd'}(\Psi)$ = Context-dependent complementarity (from Appendix HOW)
    \item $\Psi = (\Psi_I, \ldots, \Psi_F)$ = Context vector (this Appendix CTW)
\end{itemize}

\subsubsection{Summary: The 4C CORE Structure}

\begin{center}
\begin{tabular}{llll}
\toprule
\textbf{CORE} & \textbf{Question} & \textbf{Appendix} & \textbf{Output} \\
\midrule
WHO & Who has utility? & AAA & Aggregation level $L$ \\
WHAT & What is utility? & C & FEPSDE dimensions $d$ \\
HOW & How do they interact? & B & Complementarity $\Gamma(\Psi)$ \\
WHEN & When does context matter? & V & Context vector $\Psi$ \\
WHERE & Where do numbers come from? & BBB & Parameters $\Theta$, $E(\theta)$ \\
\bottomrule
\end{tabular}
\end{center}

\vspace{0.5em}
\hrule
\vspace{0.5em}
\noindent\textit{Cross-references: Appendix WHO (Aggregation Levels), 
Appendix HOW (Complementarity Structure), Appendix EST (Parameter Calibration),
Appendix WAT (FEPSDE Matrix), Appendix OPS (Operationalization), 
Appendix SUN (Falsifiable Predictions), Appendix EPS (Epistemics of Deviation)}


%==============================================================================
% BACK MATTER FOR APPENDIX V
%==============================================================================

%------------------------------------------------------------------------------
% SUMMARY BOX
%------------------------------------------------------------------------------

\section{Summary}
\label{sec:V-summary}

\begin{tcolorbox}[colback=green!10!white, colframe=green!75!black, 
    title=\textbf{Appendix CTW: Complete Summary}]

\textbf{Core Contribution:} 
Operationalization of ``context'' through eight measurable dimensions.

\textbf{The 8$\Psi$ Framework:}
\[
\boldsymbol{\Psi} = (\Psi_I, \Psi_S, \Psi_C, \Psi_K, \Psi_E, \Psi_T, \Psi_M, \Psi_F) \in [0,1]^8
\]

\begin{center}
\small
\begin{tabular}{@{}cl@{}}
$\Psi_I$ & Institutional (rules, enforcement) \\
$\Psi_S$ & Social (trust, norms) \\
$\Psi_C$ & Cognitive (skills, capacity) \\
$\Psi_K$ & Informational (transparency) \\
$\Psi_E$ & Economic (development) \\
$\Psi_T$ & Temporal (stability) \\
$\Psi_M$ & Market (openness) \\
$\Psi_F$ & Factor (flexibility) \\
\end{tabular}
\end{center}

\textbf{Key Results:}
\begin{itemize}[nosep]
    \item 8 dimensions are necessary and sufficient (Axioms Ψ-COMPLETE, Ψ-NECESSARY)
    \item Average $R^2 = 55.2\%$ for treatment effect heterogeneity
    \item Different dimensions dominate for different effects
    \item Context modulates complementarity via $\delta$-coefficients
\end{itemize}

\textbf{Practical Use:}
\begin{enumerate}[nosep]
    \item Measure $\boldsymbol{\Psi}$ using standardized indicators
    \item Calculate context-adjusted $\gamma(\Psi)$ using $\delta$-coefficients
    \item Predict intervention effectiveness based on $\Psi$-profile
    \item Target interventions where binding constraints exist
\end{enumerate}

\end{tcolorbox}

%------------------------------------------------------------------------------
% GLOSSARY OF SYMBOLS
%------------------------------------------------------------------------------

\section{Glossary of Symbols}
\label{sec:V-glossary}

\begin{longtable}{@{}clp{7cm}@{}}
\toprule
\textbf{Symbol} & \textbf{Name} & \textbf{Definition} \\
\midrule
\endhead
$\boldsymbol{\Psi}$ & Context vector & 8-dimensional vector characterizing economic environment \\
$\Psi_I$ & Institutional dimension & Rule of law, property rights, enforcement quality \\
$\Psi_S$ & Social dimension & Generalized trust, cooperative norms, social capital \\
$\Psi_C$ & Cognitive dimension & Population cognitive skills, human capital \\
$\Psi_K$ & Informational dimension & Transparency, information quality, press freedom \\
$\Psi_E$ & Economic dimension & Development level, GDP per capita, resources \\
$\Psi_T$ & Temporal dimension & Policy stability, predictability, time horizons \\
$\Psi_M$ & Market dimension & Trade openness, global integration, market scope \\
$\Psi_F$ & Flexibility dimension & Factor mobility, labor market adaptability \\
\midrule
$\delta_{dd',k}$ & Modulation coefficient & Effect of $\Psi_k$ on complementarity $\gamma_{dd'}$ \\
$\gamma(\Psi)$ & Context-adjusted complementarity & Complementarity given context $\Psi$ \\
$\bar{\gamma}$ & Baseline complementarity & Complementarity at average context ($\Psi_k = 0.5$) \\
\midrule
$\tau$ & Treatment effect & Causal effect of intervention \\
$\tau(\Psi)$ & Context-conditional effect & Treatment effect given context \\
\bottomrule
\end{longtable}

%------------------------------------------------------------------------------
% CRITICAL FOUNDATIONS (FOR CORE APPENDIX)
%------------------------------------------------------------------------------

\section{Critical Foundations}
\label{sec:V-foundations}

As a CORE appendix, Appendix CTW establishes foundations that other appendices build upon:

\begin{enumerate}
    \item \textbf{Context is Systematic, Not Residual}
    
    Treatment effect heterogeneity is not random noise---it follows predictable 
    patterns based on measurable context dimensions. This enables cross-context 
    prediction and intervention targeting.
    
    \item \textbf{Eight Dimensions Are Sufficient}
    
    The claim of sufficiency is strong but empirically supported. No ninth 
    dimension systematically improved predictions across all validation tests. 
    This does not preclude domain-specific additions for specialized applications.
    
    \item \textbf{Measurement Is Possible}
    
    Each dimension can be operationalized using existing data sources (WGI, WVS, 
    PISA, etc.). The framework is practical, not just theoretical.
    
    \item \textbf{Context Modulates Complementarity}
    
    The $\delta$-coefficients provide the link between context (Appendix CTW) and 
    complementarity (Appendix HOW). This enables the full EBF model to adjust 
    predictions across contexts.
    
    \item \textbf{Different Contexts Require Different Interventions}
    
    The same intervention (e.g., information provision) will have different effects 
    in different contexts. This is the core practical implication of the framework.
\end{enumerate}

%------------------------------------------------------------------------------
% BEHAVIORAL GAME THEORY CONNECTION
%------------------------------------------------------------------------------

\subsection{Design $>$ Equilibrium: The Mauersberger/Nagel Foundation}
\label{sec:V-design-equilibrium}

The 8$\Psi$ framework has a direct foundation in behavioral game theory.
Mauersberger \& Nagel (2018) establish a key empirical finding that validates
EBF's emphasis on context:

\begin{tcolorbox}[colback=yellow!5!white, colframe=yellow!75!black,
    title=\textbf{Design $>$ Equilibrium: Context Shifts Behavior Without Changing Payoffs}]

\textbf{Key Finding (Mauersberger \& Nagel 2018):}
\begin{quote}
``Anchors, signals, and framing shift behavior massively---often \textbf{without changing
the equilibrium}.''
\end{quote}

\textbf{Implication for $\Psi$:}
Context dimensions don't just modulate \textit{how much} behavior changes---they determine
\textit{whether} behavior changes at all, even when payoffs remain constant.

\begin{center}
\renewcommand{\arraystretch}{1.2}
\begin{tabular}{@{}p{3.5cm}p{4.5cm}p{4.5cm}@{}}
\toprule
\textbf{M\&N Concept} & \textbf{$\Psi$-Dimension} & \textbf{Mechanism} \\
\midrule
Anchor (Level-0 belief) & $\Psi_K$ (Information) & Shapes initial expectations \\
Signal quality & $\Psi_K$, $\Psi_C$ (Trust) & Credibility of coordination cues \\
Framing & $\Psi_S$ (Social norms) & Reference point for evaluation \\
Focal point visibility & $\Psi_K$, $\Psi_M$ (Market) & Salience of coordination target \\
\bottomrule
\end{tabular}
\end{center}

\end{tcolorbox}

\paragraph{Why This Matters for Interventions.}
The ``Design $>$ Equilibrium'' principle explains a puzzle in behavioral economics:

\textbf{Puzzle:} Why do small contextual changes sometimes produce large behavioral effects,
while large incentive changes sometimes produce no effect?

\textbf{Answer:} In games with strategic complementarity ($\gamma > 0$), context shifts
the \textit{anchor} (Level-0 belief), which cascades through higher reasoning levels.
The equilibrium may be unchanged, but the behavioral \textit{path} to equilibrium---and
whether agents reach it---depends critically on $\Psi$.

\begin{equation}
\text{Behavior} = f(\text{Equilibrium}, \Psi) \quad \text{not just} \quad f(\text{Equilibrium})
\tag{V-DESIGN}
\end{equation}

This is why the 8$\Psi$ framework is essential: it operationalizes the contextual factors
that determine behavior \textit{beyond} what payoff analysis alone can predict.

\textbf{Cross-Reference:} See Appendix MSC Section 11 for the full treatment of
Mauersberger \& Nagel (2018) and the worked example on Tool Adoption.

%------------------------------------------------------------------------------
% OPEN ISSUES / FUTURE WORK
%------------------------------------------------------------------------------

\section{Open Issues and Future Work}
\label{sec:V-open}

\subsection{Methodological Priorities}

\begin{enumerate}
    \item \textbf{Sub-National Measurement}
    
    Current operationalization uses country-level data. Developing regional 
    and local $\Psi$-measures would enable within-country analysis (e.g., 
    German Länder, US states, Chinese provinces).
    
    \item \textbf{Individual-Level Context}
    
    How do individual characteristics interact with aggregate $\Psi$? 
    A high-trust individual in a low-trust society faces different constraints 
    than the average.
    
    \item \textbf{Dynamic $\Psi$ Modeling}
    
    How does $\boldsymbol{\Psi}$ evolve over time? What causes institutional 
    improvement or trust decay? Endogenizing context is a major frontier.
\end{enumerate}

\subsection{Empirical Priorities}

\begin{enumerate}
    \item \textbf{More Validation Tests}
    
    Current validation covers 6 economic relationships. Extending to health, 
    education, environmental, and political outcomes would strengthen the framework.
    
    \item \textbf{$\delta$-Coefficient Estimation}
    
    Most $\delta$-coefficients are theoretically derived or LLM-MC estimated. 
    Direct econometric estimation from panel data is needed.
    
    \item \textbf{Threshold Effects}
    
    When does context ``flip'' an effect from positive to negative? 
    Identifying critical $\Psi$-thresholds would have major practical implications.
\end{enumerate}

\subsection{Theoretical Priorities}

\begin{enumerate}
    \item \textbf{Microfoundations}
    
    Why exactly do these 8 dimensions matter? Deeper theoretical grounding 
    in decision theory, game theory, and information economics would strengthen 
    the axiomatic foundations.
    
    \item \textbf{Dimension Interactions}
    
    Does $\Psi_I \times \Psi_S$ interaction matter beyond main effects? 
    Current specification assumes additivity; interactions may be important.
\end{enumerate}

%------------------------------------------------------------------------------
% REFERENCES
%------------------------------------------------------------------------------


%------------------------------------------------------------------------------
% CENTRAL REFERENCES (Required per Template)
%------------------------------------------------------------------------------

\begin{tcolorbox}[colback=blue!5!white, colframe=blue!75!black, 
    title=\textbf{Central References}]

\textbf{Glossary:} For complete symbol definitions across the EBF framework, 
see \textbf{Appendix GLS} (\textit{Glossary and Notation}).

\textbf{Bibliography:} For full reference details, see the 
\textbf{Master Bibliography} (\texttt{bcm2\_master\_references.bib}) 
in the repository.

All symbols used in this appendix are consistent with Appendix GLS definitions.

\end{tcolorbox}

\section{References for Appendix CTW}
\label{sec:V-references}

\subsection{Primary Data Sources}

\begin{itemize}[nosep]
    \item World Bank Worldwide Governance Indicators (WGI): \url{https://info.worldbank.org/governance/wgi/}
    \item World Values Survey (WVS): \url{https://www.worldvaluessurvey.org/}
    \item OECD PISA: \url{https://www.oecd.org/pisa/}
    \item Reporters Sans Frontières (RSF): \url{https://rsf.org/}
    \item Economic Policy Uncertainty (EPU): \url{https://www.policyuncertainty.com/}
    \item OECD Employment Protection: \url{https://www.oecd.org/employment/}
\end{itemize}

\subsection{Key Literature}

\begin{itemize}[nosep]
    \item Acemoglu, D., \& Robinson, J. A. (2012). \textit{Why Nations Fail}. Crown Business.
    \item Algan, Y., \& Cahuc, P. (2010). Inherited trust and growth. \textit{AER}, 100(5), 2060--2092.
    \item Bloom, N., et al. (2012). Management practices across firms and countries. \textit{AER}, 102(4), 1341--1383.
    \item Falk, A., et al. (2018). Global evidence on economic preferences. \textit{QJE}, 133(4), 1645--1692.
    \item North, D. C. (1990). \textit{Institutions, Institutional Change and Economic Performance}. Cambridge.
    \item Putnam, R. D. (1993). \textit{Making Democracy Work}. Princeton.
    \item Psacharopoulos, G., \& Patrinos, H. A. (2018). Returns to investment in education. \textit{Education Economics}, 26(5), 445--458.
\end{itemize}

\vspace{1em}
\hrule
\vspace{0.5em}
\noindent\textit{Cross-references: Appendix WHO (Aggregation Levels), 
Appendix HOW (Complementarity Structure), Appendix EST (Parameter Calibration),
Appendix WAT (FEPSDE Matrix), Appendix OPS (Operationalization), 
Appendix SUN (Falsifiable Predictions), Appendix EPS (Epistemics of Deviation)}

